\chapter{DESENVOLVIMENTO}
\label{cha:desenvolvimento}

Este capítulo descreve a arquitetura proposta e os componentes desenvolvidos para implementar o ambiente cyber-físico do Processo Tennessee Eastman sobre arquitetura OPC UA, com prova de conceito de supervisão Kubernetes. O ambiente é composto por quatro repositórios com responsabilidades bem delimitadas: \texttt{tep-plant} (simulador da planta), \texttt{tep-operator} (Operator Kubernetes), \texttt{tep-ihm} (Interface Homem-Máquina) e \texttt{tep-supervisor} (infraestrutura local).

% -------------------------------------------------------------------------
\section{Visão Geral da Arquitetura Proposta}
\label{sec:arch}

A proposta completa do ambiente TEP-CPS é estruturada sobre a hierarquia funcional da IEC~62264, cobrindo todos os cinco níveis do Modelo de Referência de Purdue, conforme o mapeamento da Figura~\ref{fig:arch_overview}:

\begin{itemize}
  \item \textbf{Nível 0 — Processo físico}: o modelo matemático do TEP (\textit{crate} \texttt{te-core}), com as equações diferenciais e o integrador RK4. Não possui nenhuma lógica de controle.
  \item \textbf{Nível 1 — Sensing e manipulação}: o \texttt{ControllerBank} com os controladores regulatórios (proporcional) e a interface gRPC que expõe medidas (\texttt{StreamMetrics}) e recebe comandos de atuação (\texttt{UpdateController}). Nível 0 e 1 coexistem dentro do repositório \texttt{tep-plant}.
  \item \textbf{Nível 2 — Controle supervisório e monitoramento}: políticas de controle declarativas aplicadas pelo Operator Kubernetes (\texttt{tep-operator}) via CRD \texttt{PLCMachine}, e interface de operação em tempo real (\texttt{tep-ihm}) para o operador humano.
  \item \textbf{Nível 3 — Gerenciamento de operações de manufatura}: fora do escopo deste trabalho.
  \item \textbf{Nível 4 — Planejamento e logística}: Kubernetes como plataforma de orquestração inteligente, responsável por executar a função objetivo composta $J$ (Equação~\ref{eq:tep_cost_extended}) e produzir as políticas supervisórias entregues ao Nível 2. Este nível é o destino arquitetural da proposta e é explorado preliminarmente no Apêndice~\ref{apendice:otimizacao}.
\end{itemize}

\begin{figure}[htbp]
\centering
\missingfigure[figwidth=0.92\textwidth]{Hierarquia funcional IEC~62264 adaptada para a proposta TEP-CPS. Nível 0: \texttt{te-core} (equações diferenciais do TEP, integrador RK4). Nível 1: \texttt{ControllerBank} (controladores P) + servidor gRPC — ambos internos ao \texttt{tep-plant}. Nível 2: \texttt{tep-operator} (Operator Kubernetes, CRD \texttt{PLCMachine}) e \texttt{tep-ihm} (IHM/dashboard WebSocket). Nível 3: fora do escopo. Nível 4: Kubernetes com função objetivo $J$ e otimização gerando políticas supervisórias. Indicar visualmente quais níveis estão implementados (0, 1, 2 — escopo deste trabalho) e quais são proposta futura (3, 4).}
\caption{Hierarquia funcional IEC~62264 do ambiente TEP-CPS: proposta completa e escopo de implementação.}
\label{fig:arch_overview}
\end{figure}

O escopo de implementação deste trabalho abrange os Níveis 0, 1 e 2. O Nível 4 — a inteligência supervisória baseada em otimização da função objetivo — é a diretriz arquitetural central da proposta e constitui a principal perspectiva de trabalho futuro.

\begin{figure}[htbp]
\centering
\missingfigure[figwidth=0.92\textwidth]{Diagrama de fluxo de dados entre os componentes implementados (Níveis 0--2). Mostrar os quatro componentes como blocos: \texttt{te-core} (Nível 0, dentro de \texttt{tep-plant}), \texttt{ControllerBank} + servidor gRPC (Nível 1, dentro de \texttt{tep-plant}), \texttt{tep-operator} (Nível 2, Kubernetes Operator), \texttt{tep-ihm} (Nível 2, IHM), e \texttt{Browser} (cliente). Setas: Nível 0 $\leftrightarrow$ Nível 1 via chamada interna \texttt{step()} e \texttt{tep\_derivatives()}; Nível 1 $\rightarrow$ \texttt{tep-ihm} via gRPC \texttt{StreamMetrics} (streaming contínuo de medidas); \texttt{tep-operator} $\rightarrow$ Nível 1 via gRPC \texttt{UpdateController} (setpoints/ganhos); \texttt{tep-ihm} $\rightarrow$ \texttt{Browser} via WebSocket (push de medidas). Destacar que \texttt{tep-operator} e \texttt{tep-ihm} não se comunicam entre si.}
\caption{Fluxo de dados entre os componentes do ambiente TEP-CPS (Níveis 0--2 implementados).}
\label{fig:data_flow}
\end{figure}

A separação entre Nível 0 e Nível 1 dentro do \texttt{tep-plant} é deliberada: \texttt{te-core} não conhece nenhum protocolo de comunicação, o que permite reutilizá-lo em contextos de otimização \textit{offline} sem carregar a dependência gRPC. No Nível 2, o \texttt{tep-operator} não conhece a \texttt{tep-ihm} e a \texttt{tep-ihm} não emite comandos de controle — a visualização não interfere no laço de controle.

% -------------------------------------------------------------------------
\section{Simulador da Planta — \texttt{tep-plant}}
\label{sec:plant}

O repositório \texttt{tep-plant} implementa o modelo matemático do TEP em Rust e expõe um servidor gRPC na porta 50051. A escolha de Rust é motivada por sua segurança de memória sem \textit{garbage collector}, desempenho próximo ao de C/C++ e sistema de tipos expressivo — características relevantes para simulação numérica de longa duração.

\subsection{Estrutura Interna}

O repositório é organizado em dois \textit{crates}:

\begin{itemize}
  \item \textbf{\texttt{te-core}}: biblioteca que implementa o modelo matemático do TEP. Contém o vetor de estado $\mathbf{y} \in \mathbb{R}^{50}$, a função de dinâmica \texttt{tep\_derivatives()} (equivalente à subrotina TEFUNC do FORTRAN original), a lógica de verificação ISD e o integrador RK4.
  \item \textbf{\texttt{tennessee-eastman-service}}: binário que instancia o modelo, executa o loop de simulação, inicia os controladores via \texttt{ControllerBank} e serve as requisições gRPC utilizando a biblioteca \textit{tonic}.
\end{itemize}

\subsection{Loop de Simulação}

O loop de simulação é executado em uma \textit{thread} dedicada. A cada iteração:

\begin{enumerate}
  \item O módulo \texttt{runtime.rs} chama \texttt{plant.step(dt)}, que internamente executa um passo RK4 com $dt = 0{,}001$~h.
  \item Após a integração, o \texttt{ControllerBank} avalia os controladores ativos via \texttt{bank.step(\&xmeas, \&mut xmv)}, calculando as novas variáveis manipuladas para o próximo passo.
  \item As medidas atualizadas são publicadas em um canal \textit{broadcast} consumido pelos clientes gRPC conectados.
  \item O estado completo pode ser serializado em um arquivo TOML (\textit{snapshot}) para reinicialização futura sem \textit{cold start}.
\end{enumerate}

O tempo simulado avança a 100× a velocidade real (1~h simulada por 36~s de relógio), controlado pela variável de ambiente \texttt{STEP\_DELAY\_MS=36}. Isso permite realizar experimentos de 20~h simuladas em aproximadamente 12~minutos de relógio.

\subsection{Controladores — \texttt{ControllerBank}}

A camada de controle é baseada na trait \texttt{Controller}:

\begin{lstlisting}[language=Rust, float=htbp, label={lst:controllerbank}, caption={Trait \texttt{Controller} e \texttt{ControllerBank} (\texttt{tep-plant})}]
pub trait Controller: Send {
    fn step(&mut self, xmeas: &[f64], xmv: &mut [f64]);
}

pub struct ControllerBank {
    controllers: Vec<Box<dyn Controller>>,
}
impl ControllerBank {
    pub fn step(&mut self, xmeas: &[f64], xmv: &mut [f64]) {
        for ctrl in &mut self.controllers {
            ctrl.step(xmeas, xmv);
        }
    }
}
\end{lstlisting}

Essa abstração separa o modelo da planta da lógica de controle, tornando o \texttt{runtime.rs} agnóstico ao tipo e quantidade de controladores. A implementação concreta \texttt{PController} realiza controle proporcional:

\begin{lstlisting}[language=Rust, float=htbp, label={lst:pcontroller}, caption={Controlador Proporcional (\texttt{PController}) --- \texttt{tep-plant}}]
pub struct PController {
    pub xmeas_idx: usize,
    pub xmv_idx:   usize,
    pub kp:        f64,
    pub setpoint:  f64,
    pub bias:      f64,
}
impl Controller for PController {
    fn step(&mut self, xmeas: &[f64], xmv: &mut [f64]) {
        let error = xmeas[self.xmeas_idx] - self.setpoint;
        xmv[self.xmv_idx] = (self.bias + self.kp * error).clamp(0.0, 100.0);
    }
}
\end{lstlisting}

No ponto de operação nominal, três malhas P são instanciadas em \texttt{main.rs}:

\begin{enumerate}
  \item \textbf{Pressão do reator} (XMEAS$_7$ → XMV$_6$): $K_p = 0{,}10$, setpoint = 2705~kPa, bias = 40,06\%.
  \item \textbf{Nível do separador} (XMEAS$_{12}$ → XMV$_7$): $K_p = 1{,}00$, setpoint = 50\%, bias = 38,10\%.
  \item \textbf{Nível do stripper} (XMEAS$_{15}$ → XMV$_8$): $K_p = 1{,}00$, setpoint = 50\%, bias = 46,50\%.
\end{enumerate}

\subsection{Servidor gRPC}

O servidor gRPC expõe quatro métodos definidos no arquivo \texttt{tep.proto}:

\begin{itemize}
  \item \texttt{StreamMetrics}: fluxo contínuo (\textit{server streaming}) de XMEAS$_{1-41}$ e XMV$_{1-12}$.
  \item \texttt{GetPlantStatus}: snapshot pontual do estado da planta.
  \item \texttt{ListControllers}: lista os controladores ativos e seus parâmetros.
  \item \texttt{UpdateController}: atualiza o setpoint ou ganho de um controlador pelo índice.
\end{itemize}

O método \texttt{UpdateController} é o canal pelo qual o operator Kubernetes age sobre a planta, fechando o laço supervisório.

\subsection{Sistema de Distúrbios}

Os distúrbios (IDV$_1$ a IDV$_{20}$) são ativados via variável de ambiente \texttt{ACTIVE\_IDV} ou dinamicamente via painel da IHM. Cada distúrbio é implementado como um modificador das condições de entrada da planta, fiel ao FORTRAN original de Downs e Vogel \cite{DOWNS1993}. A seleção dos distúrbios relevantes para o supervisor é discutida no Capítulo~\ref{cha:resultados}.

% -------------------------------------------------------------------------
\section{Operator Kubernetes — \texttt{tep-operator}}
\label{sec:operator}

O repositório \texttt{tep-operator} implementa um Operator Kubernetes em Go, construído com o framework Kubebuilder. Seu objetivo neste trabalho é atuar como \textbf{observador supervisório} da planta TEP: conecta-se à planta via gRPC, lê continuamente o estado das variáveis de processo e registra as observações no \textit{status} do CRD \texttt{PLCMachine}. A arquitetura prevê também a capacidade de intervenção via \texttt{UpdateController} — chamada gRPC que permite ajustar setpoints dos controladores — mas nos experimentos realizados o operator operou em modo de observação pura, sem emissão de comandos à planta. A camada de intervenção ativa é o horizonte natural de evolução deste componente.

\subsection{O CRD \texttt{PLCMachine}}

A \texttt{PLCMachine} é o recurso customizado que encapsula a política supervisória. Sua \textit{spec} define:

\begin{itemize}
  \item \texttt{plantAddress}: endereço gRPC da planta (e.g., \texttt{host.docker.internal:50051}).
  \item \texttt{variables}: lista de variáveis a monitorar, cada uma com nome, índice XMEAS, e faixas mínima e máxima aceitáveis.
  \item \texttt{controllers}: lista de controladores a gerenciar, com índice, setpoint nominal e ação corretiva.
\end{itemize}

Seu \textit{status} armazena:

\begin{itemize}
  \item \texttt{phase}: estado operacional (\texttt{Pending}, \texttt{Stable}, \texttt{Alarm}).
  \item \texttt{plantTime}: tempo simulado atual da planta.
  \item \texttt{lastReconcile}: timestamp do último ciclo de reconciliação.
  \item Estado observado de cada variável monitorada.
\end{itemize}

Um exemplo de manifest \texttt{PLCMachine} para o baseline TEP é apresentado no Código~\ref{lst:plcmachine}:

\begin{lstlisting}[language=yaml, float=htbp, caption={Manifest \texttt{PLCMachine} para o baseline TEP}, label={lst:plcmachine}]
apiVersion: infrastructure.tep.io/v1alpha1
kind: PLCMachine
metadata:
  name: tep-baseline
spec:
  plantAddress: "host.docker.internal:50051"
  variables:
    - name: reactor_pressure
      xmeasIndex: 7
      min: 2600.0
      max: 2900.0
    - name: separator_level
      xmeasIndex: 12
      min: 30.0
      max: 70.0
    - name: stripper_level
      xmeasIndex: 15
      min: 30.0
      max: 70.0
\end{lstlisting}

\subsection{Loop de Reconciliação}

O controller Go implementa o método \texttt{Reconcile(ctx, req)} chamado pelo \textit{controller-runtime} sempre que uma \texttt{PLCMachine} é criada ou modificada. O loop de reconciliação executa as seguintes etapas:

\begin{enumerate}
  \item \textbf{Leitura do recurso}: o controller busca a instância \texttt{PLCMachine} no API Server.
  \item \textbf{Conexão com a planta}: estabelece ou reutiliza a conexão gRPC com o endereço declarado em \texttt{spec.plantAddress}.
  \item \textbf{Observação}: chama \texttt{GetPlantStatus} para ler as 41~XMEAS e 12~XMV atuais.
  \item \textbf{Avaliação de política}: compara cada variável monitorada com seus limites declarados no CRD.
  \item \textbf{Atualização de status}: grava o estado observado — fase, tempo simulado, valores das variáveis monitoradas — em \texttt{PLCMachine.status}.
\end{enumerate}

O controller é reenfileirado automaticamente após cada ciclo, garantindo observação contínua da planta enquanto a \texttt{PLCMachine} existir no cluster.

\subsection{Conectividade com a Planta}

A planta (\texttt{tep-plant}) executa como contêiner Docker no host, enquanto o operator executa dentro do cluster Kind. Devido ao isolamento de rede do Kind, o endereço interno do Docker Compose (\texttt{te-plant:50051}) não é acessível de dentro do cluster. A solução é o endereço especial \texttt{host.docker.internal:50051}, que roteia para a interface do host Docker e alcança a planta, independentemente de onde o Kind está executando. Essa distinção é documentada como regra arquitetural do laboratório e deve ser respeitada em todos os manifests de \texttt{PLCMachine}.

% -------------------------------------------------------------------------
\section{Interface Homem-Máquina — \texttt{tep-ihm}}
\label{sec:ihm}

O repositório \texttt{tep-ihm} implementa a IHM do laboratório em Python, usando FastAPI como servidor web e WebSocket para comunicação em tempo real com o navegador.

\subsection{Arquitetura}

A IHM segue uma arquitetura assíncrona baseada em tarefas corrotinas (\textit{asyncio}):

\begin{enumerate}
  \item Uma tarefa \texttt{grpc\_reader} consome o fluxo \texttt{StreamMetrics} da planta via gRPC e publica cada mensagem em uma fila assíncrona (\texttt{asyncio.Queue}).
  \item Uma tarefa \texttt{ws\_broadcaster} retira mensagens da fila, serializa para JSON e transmite via WebSocket a todos os clientes conectados.
  \item O servidor FastAPI serve a interface HTML/JavaScript estática e gerencia as conexões WebSocket.
\end{enumerate}

Essa separação garante que o consumo gRPC seja independente do número de clientes WebSocket conectados, evitando que atrasos na rede do cliente afetem o laço de coleta de dados.

\subsection{Dashboard}

A interface web do dashboard foi desenvolvida seguindo as diretrizes da norma ISA-101 para interfaces de processo, priorizando clareza informacional sobre estética. O dashboard é composto por:

\begin{itemize}
  \item \textbf{Painel SVG da planta}: representação esquemática dos cinco subsistemas do TEP com indicação visual dos estados de alarme.
  \item \textbf{Gráficos de tendência}: séries temporais das principais variáveis (XMEAS$_7$~pressão, XMEAS$_9$~temperatura, XMEAS$_{12}$~nível separador, XMEAS$_{15}$~nível stripper) renderizadas via ECharts.
  \item \textbf{Painel de controladores}: valores atuais de cada XMV com indicação de setpoint.
  \item \textbf{Painel de distúrbios}: botões de ativação/desativação dos IDVs, com estado visual.
  \item \textbf{Painel analítico}: gráficos de composição (XMEAS$_{23-28}$) e inventários de vapor (UCVR, UCVV).
  \item \textbf{Status do PLCMachine}: exibe o estado atual do CRD via chamada ao API Server do Kubernetes (fase, última reconciliação, variáveis monitoradas).
\end{itemize}

\begin{figure}[htbp]
\centering
\missingfigure[figwidth=0.92\textwidth]{Captura de tela do dashboard da IHM em operação: painel SVG da planta (superior esquerdo) com os cinco subsistemas e indicação visual de alarmes; gráficos de tendência das variáveis XMEAS$_7$ (pressão do reator), XMEAS$_9$ (temperatura do reator), XMEAS$_{12}$ (nível do separador) e XMEAS$_{15}$ (nível do stripper) renderizados via ECharts; painel de controladores com valores atuais de XMV e setpoints; botões de ativação/desativação dos IDV$_1$--IDV$_{20}$; painel de status do PLCMachine com fase operacional, timestamp da última reconciliação e estado das variáveis monitoradas.}
\caption{Dashboard da Interface Homem-Máquina (IHM): monitoramento em tempo real do TEP.}
\label{fig:ihm_dashboard}
\end{figure}

\subsection{Registro de Dados (CSV)}

A IHM suporta gravação de todos os dados recebidos em arquivo CSV, ativada pela variável de ambiente \texttt{RECORD\_CSV=true}. O arquivo é escrito continuamente durante a sessão e pode ser exportado pelo operador via botão na interface. Os CSVs gerados são a principal fonte de dados para as análises dos experimentos descritos no Capítulo~\ref{cha:resultados}.

Cada linha do CSV contém o timestamp, o tempo simulado da planta, os valores de XMEAS$_{1-41}$ e XMV$_{1-12}$, e indicadores calculados como \texttt{deriv\_norm} e o inventário total de vapor.

% -------------------------------------------------------------------------
\section{Infraestrutura Local — \texttt{tep-supervisor}}
\label{sec:supervisor}

O repositório \texttt{tep-supervisor} reúne os arquivos de configuração de infraestrutura necessários para executar o laboratório em ambiente local Windows.

\subsection{Modo de Desenvolvimento — Docker Compose}

Para desenvolvimento e testes rápidos, a planta e a IHM são executadas via Docker Compose. O arquivo \texttt{local/docker-compose.yml} define:

\begin{itemize}
  \item Serviço \texttt{te-plant}: imagem do simulador, com portas 50051 (gRPC) e variáveis de ambiente de controle (\texttt{STEP\_DELAY\_MS}, \texttt{ACTIVE\_IDV}).
  \item Serviço \texttt{tep-ihm}: imagem da IHM, com porta 8080 (HTTP/WebSocket), variáveis \texttt{GRPC\_HOST} apontando para \texttt{te-plant:50051}, e volume para persistência dos CSVs.
\end{itemize}

A rede interna do Docker Compose permite que a IHM alcance a planta pelo nome de serviço \texttt{te-plant}, sem necessidade de configuração adicional.

\subsection{Modo de Produção — Cluster Kind}

Para execução do Operator Kubernetes, um cluster local é provisionado com \textbf{Kind} (\textit{Kubernetes in Docker}). O Kind executa um cluster Kubernetes completo dentro de contêineres Docker, sem necessidade de máquina virtual separada — adequado para desenvolvimento em Windows.

O roteiro de inicialização do laboratório com o Operator ativo é:

\begin{enumerate}
  \item Iniciar a planta e a IHM via Docker Compose.
  \item Criar o cluster Kind: \texttt{kind create cluster --name tep}.
  \item Instalar os CRDs no cluster: \texttt{kubectl apply -f config/crd/}.
  \item Construir e carregar a imagem do operator no cluster: \texttt{kind load docker-image tep-operator:latest}.
  \item Implantar o operator: \texttt{kubectl apply -f config/manager/}.
  \item Aplicar a política supervisória: \texttt{kubectl apply -f config/samples/}.
\end{enumerate}

A tabela de conectividade entre os componentes em cada modo de execução é resumida na Tabela~\ref{tab:conectividade}.

\begin{table}[ht!]
\centering
\caption{Endereços de conexão por modo de execução}
\label{tab:conectividade}
\begin{tabular}{|l|l|l|}
\hline
\textbf{Componente}       & \textbf{Docker Compose}      & \textbf{Kind + Operator} \\
\hline
IHM → planta              & \texttt{te-plant:50051}      & \texttt{te-plant:50051}  \\
Operator → planta         & N/A                          & \texttt{host.docker.internal:50051} \\
Navegador → IHM           & \texttt{localhost:8080}      & \texttt{localhost:8080}  \\
\hline
\end{tabular}
\fonte{Elaborado pelo autor.}
\end{table}

% -------------------------------------------------------------------------
\section{Considerações Finais}
\label{sec:dev_conclusao}

A arquitetura proposta realiza a separação clara entre modelo de planta, camada de controle e interface de operação, com comunicação baseada em contratos gRPC versionados. A \texttt{PLCMachine} CRD traduz o conceito de política supervisória do IEC~61499 para o paradigma declarativo do Kubernetes, e o Operator implementa o loop de reconciliação como equivalente ao runtime de implantação do padrão de Operator. O próximo capítulo apresenta os resultados dos experimentos conduzidos com essa arquitetura.
